\documentclass[UTF8,11pt]{ctexart}
\usepackage[a4paper,margin=24mm]{geometry}
\usepackage{amsmath,amssymb,amsthm,mathtools,mathrsfs}
\usepackage[all]{xy}
\usepackage{xurl,hyperref,enumitem}
\usepackage{tikz}
\usetikzlibrary{arrows.meta,cd,decorations.markings,calc,patterns}
\hypersetup{hidelinks,breaklinks=true}
\setlength{\emergencystretch}{3em}
\allowdisplaybreaks
\newtheorem*{theorem}{Theorem}
\newtheorem*{thm}{Theorem}
\newtheorem*{lemma}{Lemma}
\newtheorem*{proposition}{Proposition}
\newtheorem*{prop}{Proposition}
\newtheorem*{corollary}{Corollary}
\newtheorem*{cor}{Corollary}
\newtheorem*{claim}{Claim}
\theoremstyle{definition}
\newtheorem*{definition}{Definition}
\newtheorem*{defn}{Definition}
\newtheorem*{remark}{Remark}
\newtheorem*{example}{Example}
\newcommand{\R}{\mathbb R}
\newcommand{\C}{\mathbb C}
\newcommand{\Z}{\mathbb Z}
\newcommand{\Q}{\mathbb Q}
\newcommand{\N}{\mathbb N}
\newcommand{\F}{\mathbb F}
\newcommand{\D}{\mathbb D}
\newcommand{\bB}{\mathbb B}
\newcommand{\bC}{\mathbb C}
\newcommand{\bR}{\mathbb R}
\newcommand{\bZ}{\mathbb Z}
\newcommand{\bN}{\mathbb N}
\newcommand{\bQ}{\mathbb Q}
\newcommand{\bD}{\mathbb D}
\newcommand{\bF}{\mathbb F}
\newcommand{\CP}{\mathbb{CP}}
\newcommand{\RP}{\mathbb{RP}}
\newcommand{\sO}{\mathscr O}
\newcommand{\sI}{\mathscr I}
\newcommand{\sF}{\mathscr F}
\newcommand{\sG}{\mathscr G}
\newcommand{\sH}{\mathscr H}
\newcommand{\sR}{\mathscr R}
\newcommand{\sM}{\mathscr M}
\newcommand{\sV}{\mathscr V}
\newcommand{\sU}{\mathscr U}
\DeclareMathOperator{\Hom}{Hom}
\DeclareMathOperator{\Ext}{Ext}
\DeclareMathOperator{\Tor}{Tor}
\DeclareMathOperator{\Spec}{Spec}
\DeclareMathOperator{\Proj}{Proj}
\DeclareMathOperator{\supp}{supp}
\DeclareMathOperator{\im}{im}
\DeclareMathOperator{\id}{id}
\DeclareMathOperator{\Tr}{Tr}
\DeclareMathOperator{\tr}{tr}
\DeclareMathOperator{\rank}{rank}
\DeclareMathOperator{\ord}{ord}
\DeclareMathOperator{\dist}{dist}
\DeclareMathOperator{\End}{End}
\DeclareMathOperator{\Aut}{Aut}
\DeclareMathOperator{\Gal}{Gal}
\DeclareMathOperator{\vol}{vol}
\DeclareMathOperator{\Cl}{Cl}
\DeclareMathOperator{\coker}{coker}
\DeclareMathOperator{\codim}{codim}
\DeclareMathOperator{\divergence}{div}
\DeclareMathOperator{\Ric}{Ric}
\DeclareMathOperator{\grad}{grad}
\DeclareMathOperator{\Hess}{Hess}
\newcommand{\abs}[1]{\left\lvert#1\right\rvert}
\newcommand{\sabs}[1]{\lvert#1\rvert}
\newcommand{\babs}[1]{\bigl\lvert#1\bigr\rvert}
\newcommand{\Babs}[1]{\Bigl\lvert#1\Bigr\rvert}
\newcommand{\bbabs}[1]{\biggl\lvert#1\biggr\rvert}
\newcommand{\BBabs}[1]{\Biggl\lvert#1\Biggr\rvert}
\newcommand{\norm}[1]{\left\lVert#1\right\rVert}
\newcommand{\snorm}[1]{\lVert#1\rVert}
\newcommand{\bnorm}[1]{\bigl\lVert#1\bigr\rVert}
\newcommand{\Bnorm}[1]{\Bigl\lVert#1\Bigr\rVert}
\newcommand{\bbnorm}[1]{\biggl\lVert#1\biggr\rVert}
\newcommand{\BBnorm}[1]{\Biggl\lVert#1\Biggr\rVert}
\newcommand{\widebar}[1]{\overline{#1}}
\newcommand{\nicefrac}[2]{\frac{#1}{#2}}
\newcommand{\linnprod}[2]{\langle#1,#2\rangle}
\newcommand{\myindex}[1]{#1}
\newcommand{\glsadd}[1]{}
\newcommand{\avoidbreak}{}
\newcommand{\sourceref}[1]{\textnormal{[原文：\nolinkurl{#1}]}}
\newcommand{\thmref}[1]{\sourceref{#1}}
\newcommand{\lemmaref}[1]{\sourceref{#1}}
\newcommand{\propref}[1]{\sourceref{#1}}
\newcommand{\corref}[1]{\sourceref{#1}}
\newcommand{\defnref}[1]{\sourceref{#1}}
\newcommand{\exampleref}[1]{\sourceref{#1}}
\newcommand{\exerciseref}[1]{\sourceref{#1}}
\newcommand{\figureref}[1]{\sourceref{#1}}
\newcommand{\sectionref}[1]{\sourceref{#1}}
\newcommand{\appendixref}[1]{\sourceref{#1}}
\newcommand{\stacksref}[2]{\href{https://stacks.math.columbia.edu/tag/#1}{#2 [#1]}}

\begin{document}
\section*{035\quad Bochner--Martinelli多变量积分表示}
\subsection*{来源与计数目标}
Ji\v{r}\'i Lebl, \emph{Tasty Bits of Several Complex Variables}。从作者公开的原生 TeX 正文摘录；获取日期：2026-09-28。使用实际访问的在线正文，不把工作分支冒称冻结的印刷版。
\par\url{https://raw.githubusercontent.com/jirilebl/scv/master/scv.tex}
\par\url{https://www.jirka.org/scv/}
原文位置：Chapter 5，\emph{The Bochner--Martinelli kernel}，\texttt{thm:bochnermartinelli} 及完整证明。
\par\textbf{本文件只计一个问题：}光滑函数的 Bochner--Martinelli 积分表示（包含全纯情形），不是只引用已知核公式。
\subsection*{作者命题与人类证明}

First, let us define the Bochner--Martinelli kernel:
\begin{multline*}
\omega(\zeta,z):=\frac{(n-1)!}{(2\pi i)^n}\sum_{k=1}^n\frac{\bar\zeta_k-\bar z_k}{\|\zeta-z\|^{2n}}\\
{}\cdot d\bar\zeta_1\wedge d\zeta_1\wedge\cdots\wedge\widehat{d\bar\zeta_k}\wedge d\zeta_k\wedge\cdots\wedge d\bar\zeta_n\wedge d\zeta_n.
\end{multline*}
The notation $\widehat{d\bar\zeta_k}$ means that this term is simply left out.
\begin{thm}[Bochner--Martinelli]
Let $U\subset\C^n$ be a bounded open set with smooth boundary and let $f:\widebar U\to\C$ be a smooth function. Then for $z\in U$,
\[f(z)=\int_{\partial U}f(\zeta)\omega(\zeta,z)-\int_U\bar\partial f(\zeta)\wedge\omega(\zeta,z).\]
In particular, if $f\in\sO(U)$, then
\[f(z)=\int_{\partial U}f(\zeta)\omega(\zeta,z).\]
\end{thm}
Recall that if $\zeta=x+iy$ are the coordinates in $\C^n$, the orientation that we assigned to $\C^n$ in this book is the one corresponding to the volume form
\[dV=dx_1\wedge dy_1\wedge dx_2\wedge dy_2\wedge\cdots\wedge dx_n\wedge dy_n.\]
With this orientation,
\[d\zeta_1\wedge d\bar\zeta_1\wedge\cdots\wedge d\zeta_n\wedge d\bar\zeta_n=(-2i)^n dV,\]
and hence
\[d\bar\zeta_1\wedge d\zeta_1\wedge\cdots\wedge d\bar\zeta_n\wedge d\zeta_n=(2i)^n dV.\]
Recall the definition of $\partial$ and $\bar\partial$ from \defnref{defn:danddbarform}, and recall that $d\eta=\partial\eta+\bar\partial\eta$.
\begin{proof}
The structure of the proof is essentially the same as that of the Cauchy--Pompeiu theorem for $n=1$, although some of the formulas are more involved. Let $z\in U$ be fixed. Suppose $r>0$ is small enough so that $\overline{B_r(z)}\subset U$. Orient both $\partial U$ and $\partial B_r(z)$ positively. As $f(\zeta)\omega(\zeta,z)$ contains all the holomorphic $d\zeta_k$,
\begin{align*}
d\bigl(f(\zeta)\omega(\zeta,z)\bigr)&=\bar\partial\bigl(f(\zeta)\omega(\zeta,z)\bigr)\\
&=\bar\partial f(\zeta)\wedge\omega(\zeta,z)\\
&\quad+f(\zeta)\frac{(n-1)!}{(2\pi i)^n}\sum_{k=1}^n\frac{\partial}{\partial\bar\zeta_k}\left[\frac{\bar\zeta_k-\bar z_k}{\|\zeta-z\|^{2n}}\right]\\
&\qquad\qquad\cdot d\bar\zeta_1\wedge d\zeta_1\wedge\cdots\wedge d\bar\zeta_n\wedge d\zeta_n.
\end{align*}
We compute
\[
\sum_{k=1}^n\frac{\partial}{\partial\bar\zeta_k}\left[\frac{\bar\zeta_k-\bar z_k}{\|\zeta-z\|^{2n}}\right]=\sum_{k=1}^n\left(\frac1{\|\zeta-z\|^{2n}}-n\frac{|\zeta_k-z_k|^2}{\|\zeta-z\|^{2n+2}}\right)=0.
\]
Therefore, $d\bigl(f(\zeta)\omega(\zeta,z)\bigr)=\bar\partial f(\zeta)\wedge\omega(\zeta,z)$. We apply Stokes:
\begin{align*}
\int_{\partial U}f(\zeta)\omega(\zeta,z)-\int_{\partial B_r(z)}f(\zeta)\omega(\zeta,z)
&=\int_{U\setminus\overline{B_r(z)}}d\bigl(f(\zeta)\omega(\zeta,z)\bigr)\\
&=\int_{U\setminus\overline{B_r(z)}}\bar\partial f(\zeta)\wedge\omega(\zeta,z).
\end{align*}
Again, due to the integrability, which you showed in an exercise above, the right-hand side converges to the integral over $U$ as $r\to0$. Just as for the Cauchy--Pompeiu formula, we now need to show that the integral over $\partial B_r(z)$ goes to $f(z)$ as $r\to0$. So
\[
\int_{\partial B_r(z)}f(\zeta)\omega(\zeta,z)=f(z)\int_{\partial B_r(z)}\omega(\zeta,z)+\int_{\partial B_r(z)}\bigl(f(\zeta)-f(z)\bigr)\omega(\zeta,z).
\]
To finish the proof, we will show that $\int_{\partial B_r(z)}\omega(\zeta,z)=1$, and that the second term goes to zero. We apply Stokes again and note that the volume of $B_r(z)$ is $\frac{\pi^n}{n!}r^{2n}$.
\begin{align*}
\int_{\partial B_r(z)}\omega(\zeta,z)
&=\int_{\partial B_r(z)}\frac{(n-1)!}{(2\pi i)^n}\sum_{k=1}^n\frac{\bar\zeta_k-\bar z_k}{\|\zeta-z\|^{2n}}\\
&\qquad\cdot d\bar\zeta_1\wedge d\zeta_1\wedge\cdots\wedge\widehat{d\bar\zeta_k}\wedge d\zeta_k\wedge\cdots\wedge d\bar\zeta_n\wedge d\zeta_n\\
&=\frac{(n-1)!}{(2\pi i)^n}\frac1{r^{2n}}\int_{\partial B_r(z)}\sum_{k=1}^n(\bar\zeta_k-\bar z_k)\\
&\qquad\cdot d\bar\zeta_1\wedge d\zeta_1\wedge\cdots\wedge\widehat{d\bar\zeta_k}\wedge d\zeta_k\wedge\cdots\wedge d\bar\zeta_n\wedge d\zeta_n\\
&=\frac{(n-1)!}{(2\pi i)^n}\frac1{r^{2n}}\int_{B_r(z)}d\Bigl(\sum_{k=1}^n(\bar\zeta_k-\bar z_k)\\
&\qquad\cdot d\bar\zeta_1\wedge d\zeta_1\wedge\cdots\wedge\widehat{d\bar\zeta_k}\wedge d\zeta_k\wedge\cdots\wedge d\bar\zeta_n\wedge d\zeta_n\Bigr)\\
&=\frac{(n-1)!}{(2\pi i)^n}\frac1{r^{2n}}\int_{B_r(z)}n\,d\bar\zeta_1\wedge d\zeta_1\wedge\cdots\wedge d\bar\zeta_n\wedge d\zeta_n\\
&=\frac{(n-1)!}{(2\pi i)^n}\frac1{r^{2n}}\int_{B_r(z)}n(2i)^n\,dV=1.
\end{align*}
Next, we tackle the second term. Via the same computation as above we find
\begin{align*}
&\int_{\partial B_r(z)}\bigl(f(\zeta)-f(z)\bigr)\omega(\zeta,z)\\
&=\frac{(n-1)!}{(2\pi i)^n}\frac1{r^{2n}}\Biggl(\int_{B_r(z)}\bigl(f(\zeta)-f(z)\bigr)n\,d\bar\zeta_1\wedge d\zeta_1\wedge\cdots\wedge d\bar\zeta_n\wedge d\zeta_n\\
&\qquad+\int_{B_r(z)}\sum_{k=1}^n\frac{\partial f}{\partial\bar\zeta_k}(\zeta)(\bar\zeta_k-\bar z_k)\,d\bar\zeta_1\wedge d\zeta_1\wedge\cdots\wedge d\bar\zeta_n\wedge d\zeta_n\Biggr).
\end{align*}
As $U$ is bounded, $|f(\zeta)-f(z)|\leq M\|\zeta-z\|$ and $\left|\frac{\partial f}{\partial\bar\zeta_k}(\zeta)(\bar\zeta_k-\bar z_k)\right|\leq M\|\zeta-z\|$ for some $M$. So for all $\zeta\in B_r(z)$, we have $|f(\zeta)-f(z)|\leq Mr$ and $\left|\frac{\partial f}{\partial\bar\zeta_k}(\zeta)(\bar\zeta_k-\bar z_k)\right|\leq Mr$. Hence
\begin{multline*}
\left|\int_{\partial B_r(z)}\bigl(f(\zeta)-f(z)\bigr)\omega(\zeta,z)\right|\\
\leq\frac{(n-1)!}{(2\pi)^n}\frac1{r^{2n}}\left(\int_{B_r(z)}n2^nMr\,dV+\int_{B_r(z)}n2^nMr\,dV\right)=2Mr.
\end{multline*}
Therefore, this term goes to zero as $r\to0$.
\end{proof}

\subsection*{整理说明（非作者正文）}
保留核定义、体积形式的取向约定、闭性计算、两次 Stokes 应用、球面归一化常数与误差项趋零，未省略归一化推导。标准先修为微分形式的外微分、Stokes 定理、球体体积公式、光滑函数局部 Lipschitz 估计，以及阶为 $2n-1$ 的局部奇核在 $2n$ 实维下的可积性；原书将最后一项列作习题，本条不声称重证。原文取向脚注的历史评论未取入，实际取向及正负号全部保留。

难度依据是选择带有正确形式次数的奇异核并把局部留数式的归一化和余项控制转成全局表示；不是对现成积分公式的直接数值应用。正文仅将长的楔积表达式折行，完整保留每一等式与推导次序。
\subsection*{可修改的简短标注（非作者正文）}
$c$：多复变量函数的边界与内部积分表示。

$a$：微分形式；$\partial/\bar\partial$ 分解；Stokes 定理；奇异积分核；球面归一化；局部误差估计。

\texttt{cross\_theory\_status = yes}。把复解析再生性质编码为微分形式核的闭性，通过去小球后的 Stokes 公式与球面极限返回函数点值。位置：核的导数求和、球面积分为1及最后的 $2Mr$ 估计。

全部标注均为非穷尽、可修改初稿；未标注不表示未使用。
\subsection*{许可与修改记录}
原数学正文作者：Ji\v{r}\'i Lebl。作者网站及源书声明允许 CC BY-SA 4.0 或 CC BY-NC-SA 4.0 双许可；本条选择 CC BY-SA 4.0，整理部分亦采用同一许可。\url{https://creativecommons.org/licenses/by-sa/4.0/}。2026-09-28：助手截取题证、调整独立排版，加入来源、范围与初稿标注；不暗示作者认可本次整理。保留的是所用 TeX 正文摘录，未将其称为整书原件。
\end{document}
